\section{Experimental Setup}
\label{sec:experimental-setup}

\subsection{Datasets}

We use five BEIR datasets \citep{thakur2021beir}: NFCorpus, SciFact, ArguAna, SCIDOCS, and FiQA-2018.
These datasets cover scientific, biomedical, argument retrieval, citation-related, and financial question-answering retrieval settings.
The current notebook uses the first 50 queries from each dataset to keep the early experiment practical on Kaggle.
Table~\ref{tab:datasets} lists the dataset identifiers used by the code.

\begin{table}[t]
  \centering
  \begin{tabular}{ll}
    \toprule
    Dataset & Code name \\
    \midrule
    NFCorpus & \texttt{nfcorpus} \\
    SciFact & \texttt{scifact} \\
    ArguAna & \texttt{arguana} \\
    SCIDOCS & \texttt{scidocs} \\
    FiQA-2018 & \texttt{fiqa} \\
    \bottomrule
  \end{tabular}
  \caption{Datasets used in the 50-query early experiment.}
  \label{tab:datasets}
\end{table}

\subsection{Model and Inference Settings}

All experiments use the public \texttt{Soyoung97/ListT5-base} checkpoint.
The input candidates are BM25 top-100 files from the BEIR-format dataset repository used by the ListT5 codebase.
The fixed ListT5 settings are listwise group size $m=5$, output size $r=2$, reranking depth 10, maximum input sequence length 512, and NDCG@10 evaluation.
The 512-token cap is applied to every dataset to keep the early run consistent and practical in the low-resource notebook setting.
For the BM25 top-k experiment, the input top-k value is changed to 40, 60, or 80 and compared with sequential top-100 as the local baseline.

\subsection{Compared Conditions}

The experiment grid has two parts.
First, the BM25 top-k comparison evaluates sequential grouping at BM25 top-40, top-60, top-80, and top-100.
Second, the grouping comparison evaluates sequential and score-balanced grouping at BM25 top-100.
Sequential top-100 is reused as the local baseline in both comparisons, so the full grid has five jobs per dataset: sequential top-40, top-60, top-80, top-100, and score-balanced top-100.
The reported runs were mainly executed from notebooks for Kaggle-style experimentation.
The repository also contains a Python codebase version of the grouping experiment, which can be used when moving the workflow from interactive runs to scripted execution.
